%! Function Model Construction and Application — Unit 1, Lesson 1.14 — Lesson Plan
\documentclass[10pt]{article}
\usepackage{precalculus-article}
\usepackage{precalculus-boxes}
\usepackage{graphicx}   % -article does NOT load graphicx; needed for thumbnails/screenshots
\graphicspath{{images/}}
\newcommand{\TallMath}[1]{$\displaystyle #1\rule[-1.4em]{0pt}{3.2em}$}
\newcommand{\CourseName}{Precalculus}
\newcommand{\SchoolYear}{2026--2027}
\newcommand{\MeetingLength}{55 minutes}
\newcommand{\UnitNumberName}{Unit 1: Polynomial and Rational Functions \quad}
\newcommand{\LessonNumberName}{Lesson 1.14: Function Model Construction and Application}

\begin{document}
\begin{center}
    {\Large\bfseries \CourseName: \SchoolYear \\
    {\small\bfseries \UnitNumberName \LessonNumberName}}
\end{center}

\begin{tcolorbox}[colback=sky, colframe=navy, boxrule=0.9pt, arc=2mm,
    left=2mm, right=2mm, top=1.5mm, bottom=1.5mm]
    \textbf{Primary Objective:} Students will construct polynomial, piecewise, and rational
    function models from contextual descriptions using restrictions and inverse-proportion
    relationships, and apply those models to predict values and compute average rates of change
    with appropriate units.
    \hfill\textit{\small Skills: MP 1.C, MP 3.B}
\end{tcolorbox}

\renewcommand{\arraystretch}{1.6}
\begin{skillbox}[Priority Ideas \& Skills]{goldbox}
\begin{tabularx}{\linewidth}{@{}X X@{}}
\textbf{AP Skills} &
\textbf{Key Understandings} \\
\midrule
\textbf{MP 1.C} --- Construct new functions using transformations, compositions, or regression. \newline
\textbf{MP 3.B} --- Apply numerical results in context using appropriate units. &
\begin{itemize}[leftmargin=1.2em,itemsep=2pt,topsep=0pt]
  \item A model is built from a contextual restriction by expressing one variable in terms
        of another, then substituting into the objective formula. (EK 1.14.A.1, A.2)
  \item Inversely proportional quantities $F \propto 1/d^n$ are modeled by rational functions
        $r(d) = k/d^n$, with domain $d > 0$. (EK 1.14.B.1)
  \item The average rate of change $\frac{f(b)-f(a)}{b-a}$ has units equal to output units
        divided by input units and describes the model's behavior over an interval. (EK 1.14.C.1)
\end{itemize} \\
\end{tabularx}
\end{skillbox}

\begin{skillbox}[{Vocabulary, Concepts, \& Theorems}]{greenbox}
\renewcommand{\arraystretch}{1.4}
\begin{tabularx}{\linewidth}{@{} >{\bfseries\color{navy}}p{0.28\linewidth} X @{}}
function model        & A function constructed to represent a real-world scenario; used to predict values. \\
restriction           & A constraint on input values derived from the context (e.g., length must be positive). \\
regression            & A technology-assisted technique for fitting a curve of a given type to a data set. \\
average rate of change & $\dfrac{f(b)-f(a)}{b-a}$; slope of the secant line; units = output/input. \\
inverse proportion    & A relationship $y = k/x^n$; as $x$ increases, $y$ decreases. \\
domain restriction    & The set of allowable inputs dictated by the contextual meaning of the variable. \\
piecewise model       & A function defined by different formulas on different intervals of its domain. \\
\end{tabularx}
\end{skillbox}

\begin{skillbox}[Activate Prior Knowledge \& Spiral Review]{sky}
    \begin{tabularx}{\linewidth}{@{}X@{}}
        Spiral review (warm-up) covers: (1) computing average rate of change from a table with
        correct units (from 1.1/1.4); (2) evaluating a polynomial at a given input (from 1.2);
        (3) recognizing inverse-square language as a rational model (bridges to 1.14.B).
    \end{tabularx}
\end{skillbox}

\begin{skillbox}[Hook]{sky}
``A farmer has 80 feet of fencing and wants the biggest pen possible. One side is against the barn
--- no fencing needed there. Can a function tell us the answer?''

Write the scenario on the board. Have students predict: would the optimal shape be a square?
Why or why not? Take 3--4 verbal responses. Then say: \textit{``Let's use a function model to
find out exactly.''}
\end{skillbox}

\begin{skillbox}[Lesson]{sky}
\begin{multicols}{2}
\textbf{Part 1 (15 min): Model from Restrictions (Notes Sec.\ 1).}

\begin{enumerate}[leftmargin=1.2em,itemsep=2pt,topsep=2pt]
  \item Label the pen: two widths $w$, one length $\ell$, barn on fourth side.
  \item Constraint: $2w + \ell = 80 \Rightarrow \ell = 80 - 2w$.
  \item Objective: $A = w\ell \Rightarrow A(w) = -2w^2 + 80w$.
  \item Domain: $0 < w < 40$.
  \item Vertex: $w = 20$ ft, $A(20) = 800$ ft$^2$.
\end{enumerate}

Pose: \textbf{``Is this a square pen?''} ($\ell = 40 \ne w = 20$, so no.)

\columnbreak

\textbf{Part 2 (10 min): Rational Model (Notes Sec.\ 2).}

Two charged objects: force $\propto 1/d^2$, so $r(d) = k/d^2$.
\begin{itemize}[leftmargin=1.2em,itemsep=2pt,topsep=2pt]
  \item Domain: $d > 0$ (physical distance must be positive).
  \item As $d \to 0^+$: force $\to \infty$. As $d \to \infty$: force $\to 0$.
  \item Complete table with $k = 900$: $d = 1, 3, 5, 10$.
\end{itemize}

\textbf{Part 3 (10 min): Apply the Model (Notes Sec.\ 3).}

Return to $A(w) = -2w^2 + 80w$:
\begin{itemize}[leftmargin=1.2em,itemsep=2pt,topsep=2pt]
  \item Evaluate $A(10) = 600$ ft$^2$ and $A(20) = 800$ ft$^2$.
  \item ARC on $[10, 20]$: $\frac{800-600}{10} = 20$ ft$^2$/ft.
  \item Interpret: each foot of additional width yields 20 ft$^2$ more area, on average, between
        $w = 10$ and $w = 20$.
\end{itemize}
\end{multicols}
\end{skillbox}

\begin{skillbox}[Active Monitoring]{redbox}
\begin{itemize}[leftmargin=1.2em,itemsep=2pt,topsep=0pt]
  \item After Part 1 --- cold call: \textbf{``Why is the domain $0 < w < 40$, not $w > 0$?''}
        (If $w \geq 40$, $\ell \leq 0$; no pen possible.)
  \item After Part 2 --- cold call: \textbf{``Why can't $d = 0$ in your rational model?''}
        (Division by zero; also physically impossible.)
  \item After Part 3 --- cold call: \textbf{``What units does the average rate of change have?''}
        (ft$^2$/ft. Push students to read units as ``square feet per foot of width.'')
  \item Watch for: students computing ARC without including units; students plugging into the
        original unsimplified expression rather than $A(w) = -2w^2 + 80w$.
\end{itemize}
\end{skillbox}

\begin{skillbox}[Group Work \& Differentiation]{redbox}
\begin{multicols}{3}
\textbf{Tier R --- Remediate}
\begin{itemize}[leftmargin=1em,itemsep=2pt,topsep=0pt]
  \item Area model pre-set up: evaluate $A(w) = -2w^2 + 60w$ at four widths.
  \item Compute ARC on $[5, 15]$.
  \item Identify the approximate maximum from the table.
\end{itemize}

\columnbreak

\textbf{Tier A --- Apply}
\begin{itemize}[leftmargin=1em,itemsep=2pt,topsep=0pt]
  \item Write rational model $C(d) = 2000/d^2$ from description.
  \item State domain restriction with explanation.
  \item Complete table; compute ARC on $[2, 4]$ with units.
\end{itemize}

\columnbreak

\textbf{Tier E --- Extend}
\begin{itemize}[leftmargin=1em,itemsep=2pt,topsep=0pt]
  \item Construct quadratic model by solving $3 \times 3$ system from data points.
  \item State modeling assumptions.
  \item Predict at $x = 4$ (extrapolation) and comment.
\end{itemize}
\end{multicols}
\end{skillbox}

\begin{skillbox}[Individual Work \& Assessment]{redbox}
\textbf{Exit Ticket} (3 questions, 1 page):
\begin{enumerate}[leftmargin=1.2em,itemsep=2pt,topsep=2pt]
  \item Evaluate $f(x) = -2x^2 + 12x$ at $x = 3$ and $x = 5$; compute ARC on $[3, 5]$ with units.
  \item State domain restriction for $P(d) = 50/d^2$ in context.
  \item Interpret a negative ARC in the pen design context.
\end{enumerate}

Collect exit tickets as students leave. Sort into three piles:
\begin{itemize}[leftmargin=1.2em,itemsep=2pt,topsep=2pt]
  \item \textbf{Missing units only} --- quick fix in warm-up next class.
  \item \textbf{Incorrect ARC computation} --- small-group pull during Lesson 1.14 practice.
  \item \textbf{All correct} --- ready for Unit 1 review.
\end{itemize}
\end{skillbox}

\begin{skillbox}[Reinforcement \& Extension]{goldbox}
\textbf{Homework (6 problems):}
\begin{itemize}[leftmargin=1.2em,itemsep=2pt,topsep=2pt]
  \item (1) Model from restrictions (two-sided garden): derive $A(w)$, domain, vertex.
  \item (2) Rational gravitational model: domain, end behavior, inverse-square ratio.
  \item (3) Projectile $s(t) = -16t^2 + 64t + 5$: evaluate, ARC, interpret zero ARC.
  \item (4) Tank draining: interpret sign of ARC; estimate output from average.
  \item (5) Piecewise parking-garage model: write, evaluate, classify.
  \item (6) AP-style MC on ARC of a revenue model.
\end{itemize}

\textbf{Preview --- Unit 1 Sample Test:} Lesson 1.14 is the final lesson of Unit 1.
After returning homework, distribute the Unit 1 Sample Test for in-class or at-home review.
Topics covered on the test span 1.1--1.14: polynomial behavior, rational functions (VA, holes,
end behavior), and function model construction and application.
\end{skillbox}

\end{document}
